\documentclass[11pt]{article}

\usepackage[margin=1in]{geometry}
\usepackage{amsmath}
\usepackage{amssymb}
\usepackage{booktabs}
\usepackage{siunitx}
\usepackage{microtype}
\usepackage[hidelinks]{hyperref}

\emergencystretch=3em

\newcommand{\ii}{\mathrm{i}}
\newcommand{\dd}{\,\mathrm{d}}
\newcommand{\Rw}{R_{\mathrm{w}}}
\DeclareMathOperator{\sech}{sech}

\title{Michell's Wave-Resistance Integral and the Free-Wave\\
Amplitude Grid from a B-spline Control Net}
\author{\texttt{michell}}
\date{}

\begin{document}
\maketitle

\begin{abstract}
This note documents the core numerical calculation performed by the
\texttt{michell} crate: given a ship hull represented as a tensor-product
B-spline half-breadth surface, it computes the thin-ship (Michell) wave
resistance and the far-field free-wave elevation pattern. The distinguishing
feature of the implementation is that the inner Michell integrals are evaluated
\emph{exactly} --- in closed form per knot span --- because a polynomial
B-spline hull makes the integrand a polynomial multiplied by an oscillatory /
exponentially decaying kernel on each span. The only numerical error lives in a
single smooth one-dimensional outer quadrature. Section references point at the
Rust source so the formulae and the code can be read side by side.
\end{abstract}

\tableofcontents

\section{Geometry, conventions, and constants}
\label{sec:conventions}

The hull is a slender body symmetric about its centerplane. Its starboard
surface is given as a single-valued \emph{half-breadth} (half-beam) function
\begin{equation}
  y = f(x, z) \ge 0,
\end{equation}
where
\begin{itemize}
  \item $x$ runs along the hull, with the ship advancing toward $+x$ (the bow is
        the high-$x$ end of the hull file; the wake trails toward $-x$);
  \item $z$ is measured \emph{downward} from the undisturbed waterline, so the
        wetted surface occupies $z \in [0, T]$ with $T$ the draft;
  \item $y$ is the transverse (port/starboard) coordinate.
\end{itemize}
The surface domain is $x \in [x_0, x_1]$ (length $L = x_1 - x_0$) and
$z \in [0, T]$. Only the polynomial (non-rational) B-spline surface of
Section~\ref{sec:bspline} is admitted; rational NURBS weights are deliberately
excluded, because it is precisely the piecewise-polynomial structure that makes
the inner integrals closed-form.

\paragraph{Operating conditions and constants.}
Table~\ref{tab:constants} lists the physical inputs. The single derived quantity
that governs the whole calculation is the \emph{fundamental wavenumber}
\begin{equation}
  \boxed{\;\nu \;=\; \frac{g}{U^{2}}\;}
  \qquad\text{(\texttt{nu} in the source),}
  \label{eq:nu}
\end{equation}
the wavenumber of the transverse ($\theta = 0$) waves of a steadily advancing
ship in deep water. Its reciprocal length scale is the transverse wavelength
$\lambda_{\mathrm{tr}} = 2\pi/\nu = 2\pi U^{2}/g$. The length Froude number
$\mathrm{Fr} = U/\sqrt{gL}$ is the usual dimensionless speed.

\begin{table}[h]
  \centering
  \begin{tabular}{llll}
    \toprule
    Symbol & Meaning & Units & Source \\
    \midrule
    $U$   & ship speed                 & \si{m/s}       & \texttt{Conditions::speed} \\
    $g$   & gravitational acceleration & \si{m/s^2}     & \texttt{9.80665} \\
    $\rho$& fluid density              & \si{kg/m^3}    & \texttt{Fluid::density} ($\approx 1025.9$ seawater) \\
    $\nu$ & fundamental wavenumber $g/U^2$ & \si{1/m}   & derived, Eq.~\eqref{eq:nu} \\
    $L$   & waterline length           & \si{m}         & $x_1 - x_0$ \\
    $T$   & draft                      & \si{m}         & $z$-domain upper bound \\
    \bottomrule
  \end{tabular}
  \caption{Physical inputs (\texttt{conditions.rs}). Gravity and density carry
  the ITTC 15\,\si{\celsius} default values.}
  \label{tab:constants}
\end{table}

\section{The B-spline control-net representation}
\label{sec:bspline}

The half-breadth surface is a clamped tensor-product B-spline of degree $p$ in
$x$ and $q$ in $z$:
\begin{equation}
  f(x, z) \;=\; \sum_{i=0}^{n_x-1}\sum_{j=0}^{n_z-1}
     P_{ij}\, N_{i,p}(x)\, N_{j,q}(z),
  \label{eq:surface}
\end{equation}
where $P_{ij} \in \mathbb{R}$ is the \emph{control net} of half-beam values
(stored row-major with the $z$ index fastest, \texttt{control[i*n\_ctrl\_z + j]}),
and $N_{i,p}$, $N_{j,q}$ are the B-spline basis functions built from the knot
vectors $\{\xi_k\}$ (in $x$) and $\{\zeta_k\}$ (in $z$). Both knot vectors are
\emph{clamped}: the first and last knots each repeat $p+1$ (resp.\ $q+1$) times,
so the surface interpolates the corner control points and the parametric
coordinate coincides with the physical coordinate. An interior knot of
multiplicity up to $p$ is allowed --- multiplicity $p$ drops the surface to
$C^{0}$ and is exactly how a \emph{chine} (hard crease) is represented.

Because the control net is non-negative (a sufficient condition for
$f \ge 0$ by the convex-hull property) and $z_0 = 0$ (the domain starts at the
waterline), the object is a valid hull. Basis functions and their derivatives
are evaluated with the standard Piegl--Tiller recurrences
(\texttt{bspline.rs}, \texttt{ders\_basis}).

\subsection{From control net to per-span polynomial}
\label{sec:spanpoly}

The knot vectors partition the domain into a grid of non-empty rectangular
\emph{knot spans}. On a single span
$[\xi_{s}, \xi_{s+1}] \times [\zeta_{t}, \zeta_{t+1}]$ the surface
\eqref{eq:surface} is an ordinary bivariate polynomial of degree $(p, q)$. The
crate materialises that polynomial once, from the mixed partial derivatives
evaluated at the span's lower-left corner $(x_0, z_0) = (\xi_s, \zeta_t)$:
\begin{equation}
  D_{ab} \;=\; \left.\frac{\partial^{\,a+b} f}{\partial x^{a}\,\partial z^{b}}
    \right|_{(x_0, z_0)},
  \qquad 0 \le a \le p,\; 0 \le b \le q,
  \label{eq:corner}
\end{equation}
(\texttt{BSplineSurface::corner\_partials}). Taylor's theorem then reconstructs
the surface exactly on the span,
\begin{equation}
  f(x, z) \;=\; \sum_{a=0}^{p}\sum_{b=0}^{q}
    \frac{D_{ab}}{a!\,b!}\,(x - x_0)^{a}(z - z_0)^{b}.
  \label{eq:taylor}
\end{equation}

Michell's integral needs the longitudinal slope $\partial f/\partial x$.
Differentiating \eqref{eq:taylor} and reindexing gives, on each span, the
polynomial
\begin{equation}
  \boxed{\;
  f_x(x, z) \;\equiv\; \frac{\partial f}{\partial x}
    \;=\; \sum_{a=0}^{p-1}\sum_{b=0}^{q}
      c_{ab}\,(x - x_0)^{a}(z - z_0)^{b},
  \qquad
  c_{ab} = \frac{D_{a+1,\,b}}{a!\,b!}. \;}
  \label{eq:fxcoeff}
\end{equation}
These coefficients $c_{ab}$ are precomputed for every span pair and stored
flattened as \texttt{fx\_coeff} (\texttt{hull.rs}), indexed
$[((s_x\, n_z + s_z)\,p + a)(q+1) + b]$. This array is the sole geometric input
to the wave calculation: the entire hull enters Michell's integral through the
per-span slope polynomials \eqref{eq:fxcoeff}.

\section{Michell's wave-resistance integral}
\label{sec:michell}

\subsection{The classical form}

For a thin ship in steady forward motion in calm, deep water, Michell's integral
gives the wave-making resistance as
\begin{equation}
  \Rw \;=\; \frac{4\rho g^{2}}{\pi U^{2}}
    \int_{1}^{\infty} \bigl(I(\lambda)^{2} + J(\lambda)^{2}\bigr)\,
      \frac{\lambda^{2}}{\sqrt{\lambda^{2}-1}}\dd\lambda,
  \label{eq:michell-lambda}
\end{equation}
where the real functions $I$ and $J$ are the cosine and sine parts of the
complex \emph{free-wave amplitude function}
\begin{equation}
  I(\lambda) + \ii\,J(\lambda)
    \;=\; \iint_{\text{hull}}
      \frac{\partial f}{\partial x}\,
      e^{-\nu\lambda^{2} z}\,
      e^{\,\ii\,\nu\lambda x}
      \dd x \dd z.
  \label{eq:innerdef}
\end{equation}
Here the integration variable $\lambda \ge 1$ is the secant of the wave
propagation angle $\theta$; each $\lambda$ selects one plane-wave component of
the ship's Kelvin pattern. The kernel of \eqref{eq:innerdef} has the two
signatures of linear ship-wave theory:
\begin{itemize}
  \item a horizontal oscillation $e^{\,\ii\,\nu\lambda x}$ with wavenumber
        $k_x = \nu\lambda$ along the track, and
  \item a vertical decay $e^{-\nu\lambda^{2} z}$ with rate
        $\kappa = \nu\lambda^{2}$, the deep-water statement that shorter,
        more-oblique waves are confined nearer the surface.
\end{itemize}
(See \texttt{michell.rs} module docs; cf.\ Tuck 1989; Dambrine, Pierre \&
Rousseaux 2016.)

\subsection{Secant substitution}

The integrand of \eqref{eq:michell-lambda} has an integrable singularity at
$\lambda = 1$. Substituting $\lambda = \sec\theta$
($\dd\lambda = \sec\theta\tan\theta\dd\theta$,
$\sqrt{\lambda^2-1} = \tan\theta$) sends
$\lambda^2/\sqrt{\lambda^2-1}\dd\lambda \to \sec^{3}\theta\dd\theta$ and removes
the singularity, giving the form the code actually integrates:
\begin{equation}
  \boxed{\;
  \Rw \;=\; \frac{4\rho g^{2}}{\pi U^{2}}
    \int_{0}^{\pi/2}
      \bigl|F(\sec\theta)\bigr|^{2}\,\sec^{3}\theta \dd\theta,
  \qquad
  F(\lambda) \equiv I(\lambda) + \ii\,J(\lambda). \;}
  \label{eq:michell-theta}
\end{equation}
The prefactor $4\rho g^{2}/(\pi U^{2})$ is \texttt{coeff} in
\texttt{multihull\_wave\_resistance\_with}; the integral is accumulated in
\texttt{integrate\_outer}. This outer integral is smooth, and is the
\emph{only} place numerical quadrature error enters.

\subsection{Exact evaluation of the inner integral}
\label{sec:inner}

Substitute the per-span slope polynomial \eqref{eq:fxcoeff} into the definition
\eqref{eq:innerdef}. Because $f_x$ is piecewise polynomial and the kernel
separates in $x$ and $z$, the double integral factorises span-by-span into a
sum of products of one-dimensional \emph{moment integrals}. With
$k_x = \nu\lambda$ and $\kappa = \nu\lambda^{2}$, define on a span of lengths
$h_x, h_z$ (local coordinates $X = x-x_0 \in [0,h_x]$, $Z = z-z_0 \in [0,h_z]$)
the oscillatory and exponential moments
\begin{align}
  M_a(k_x, h_x) &= \int_{0}^{h_x} X^{a}\, e^{\,\ii k_x X}\dd X,
    &  a &= 0,\dots,p-1, \label{eq:Mdef}\\
  N_b(\kappa, h_z) &= \int_{0}^{h_z} Z^{b}\, e^{-\kappa Z}\dd Z,
    &  b &= 0,\dots,q. \label{eq:Ndef}
\end{align}
Then the contribution of a single hull to the amplitude function is the exact
finite sum
\begin{equation}
  \boxed{\;
  F(\lambda) \;=\;
    \sum_{s\in\text{spans}_x}
      e^{\,\ii k_x (\xi_{s} - x_c)}
      \sum_{a=0}^{p-1} M_a(k_x, h_x^{(s)})
      \underbrace{\sum_{t\in\text{spans}_z}\sum_{b=0}^{q}
        c_{ab}^{(s,t)}\;
        e^{-\kappa\,\zeta_{t}}\, N_b(\kappa, h_z^{(t)})}_{\textstyle g_a^{(s)}}.
  \;}
  \label{eq:Fexact}
\end{equation}
This is exactly the double loop in \texttt{InnerIntegral::eval}: the inner
bracket $g_a^{(s)}$ contracts the slope coefficients $c_{ab}$ against the
decay-shifted $z$-moments $e^{-\kappa\zeta_t} N_b$ (precomputed once per
$\lambda$ across all $z$-spans, array \texttt{zm}); the outer sum weights them by
the $x$-moments $M_a$ and the per-span phase factor. The reference phase
$x_c = \tfrac12(x_0 + x_1)$ is the hull midpoint; shifting it multiplies $F$ by a
unit-modulus factor and leaves $|F|^{2}$ --- hence $\Rw$ --- invariant.

The upshot: for a polynomial hull, $I$ and $J$ carry \emph{no} discretisation
error. Every span integral is a closed form (Section~\ref{sec:moments}). The
number of inner evaluations equals (\#$x$-spans)$\times$(\#$z$-spans)$\times
p\times(q{+}1)$ complex operations per $\lambda$.

\subsection{Closed-form span moments}
\label{sec:moments}

The moments \eqref{eq:Mdef}--\eqref{eq:Ndef} are computed in \texttt{moments.rs}.
Each has a stable large-argument recurrence obtained by integration by parts,
and a small-argument power series used to avoid catastrophic cancellation
(switch-over at $|k_x h_x|,\,\kappa h_z = 4$).

\paragraph{Oscillatory moments.} Integration by parts of \eqref{eq:Mdef} gives
the upward recurrence
\begin{equation}
  M_0 = \frac{e^{\,\ii k_x h_x} - 1}{\ii k_x},
  \qquad
  M_a = \frac{h_x^{a}\,e^{\,\ii k_x h_x} - a\,M_{a-1}}{\ii k_x},
  \label{eq:Mrec}
\end{equation}
and for small $|k_x h_x|$ the series
\begin{equation}
  M_a = h_x^{a+1}\sum_{m\ge 0}
    \frac{(\ii k_x h_x)^{m}}{m!\,(a+m+1)}.
  \label{eq:Mseries}
\end{equation}

\paragraph{Exponential moments.} Likewise for \eqref{eq:Ndef} with $\kappa \ge 0$,
\begin{equation}
  N_0 = \frac{1 - e^{-\kappa h_z}}{\kappa},
  \qquad
  N_b = \frac{b\,N_{b-1} - h_z^{b}\,e^{-\kappa h_z}}{\kappa},
  \label{eq:Nrec}
\end{equation}
and for small $\kappa h_z$,
\begin{equation}
  N_b = h_z^{b+1}\sum_{m\ge 0}
    \frac{(-\kappa h_z)^{m}}{m!\,(b+m+1)}.
  \label{eq:Nseries}
\end{equation}
As $\kappa \to \infty$ (deeply oblique waves) $N_b \to b!/\kappa^{b+1} \to 0$,
so deep parts of the hull decouple gracefully; the code short-circuits a span
once $e^{-\kappa\zeta_t}$ underflows.

\subsection{Outer quadrature}
\label{sec:outer}

The smooth $\theta$-integral \eqref{eq:michell-theta} is evaluated by marching
adaptive Gauss--Legendre panels (16 nodes each). The panel width is set to a
fixed fraction of one local oscillation period of the integrand,
\begin{equation}
  \Delta\theta \;\lesssim\; \frac{2\pi}{r(\theta)},
  \qquad
  r(\theta) = 2\nu\sec\theta\tan\theta\,(x_{1/2} + T\sec\theta) + \dots,
\end{equation}
where $r(\theta)$ is the analytic phase rate of $|F|^{2}\sec^{3}\theta$
(dominated near the track by the longitudinal bandwidth $x_{1/2}$ and the decay
depth $T$). The pass is repeated with the panels halved until the change between
successive passes drops below the requested relative tolerance
(\texttt{rel\_tol}, default $10^{-5}$). The tail in $\lambda = \sec\theta$ is
truncated once a full multi-period window of the oscillatory integrand
contributes negligibly, giving the reported \texttt{max\_lambda}.

\subsection{Fleets and interference}
\label{sec:multihull}

For a multihull / fleet, the far-field free-wave amplitudes superpose linearly.
A member hull $j$ at longitudinal offset $\Delta x_j$ and transverse position
$y_j$ contributes
$F_j(\lambda)\,\exp[\ii\nu(\lambda\,\Delta x_j \pm \lambda\sqrt{\lambda^2-1}\,y_j)]$
to the wave system propagating at $\pm\theta$ off the track. Writing
$k_x = \nu\lambda$ and $k_y = \nu\lambda\sqrt{\lambda^2-1}$, the resistance
integrand replaces $|F|^2$ in \eqref{eq:michell-theta} by the mean of the two
signed systems,
\begin{equation}
  |A(\lambda)|^{2} \;=\;
    \tfrac12\Bigl(
      \bigl|\textstyle\sum_j F_j\, e^{\,\ii(k_x\Delta x_j + k_y y_j)}\bigr|^{2}
    + \bigl|\textstyle\sum_j F_j\, e^{\,\ii(k_x\Delta x_j - k_y y_j)}\bigr|^{2}
    \Bigr),
  \label{eq:multi}
\end{equation}
which is \texttt{amp\_sq} in the source. The two halves differ only for fleets
that are not mirror-symmetric about their mean centerplane (e.g.\ staggered
pairs). For two identical hulls a transverse separation $s$ apart this reduces to
the classical catamaran interference factor
$4\cos^{2}\!\bigl(\tfrac12\nu s\,\lambda\sqrt{\lambda^2-1}\bigr)$.

\section{The free-wave amplitude and elevation grid}
\label{sec:spectrum}

The same amplitude functions $F_j$ that produce the resistance also produce the
visible wave pattern (the "wake"). This is handled in \texttt{spectrum.rs}.

\subsection{Free-wave spectrum}

Far behind the hull the disturbance is a superposition of free plane waves, one
per propagation angle $\theta \in (-\tfrac\pi2, \tfrac\pi2)$ measured from the
track. Deep-water dispersion pins each component's wavenumber and wavevector at
\begin{equation}
  k(\theta) = \nu\sec^{2}\theta,
  \qquad
  \mathbf{k}(\theta) = \nu\sec\theta\,(1, \tan\theta),
\end{equation}
and the complex amplitude \emph{density} is
\begin{equation}
  \boxed{\;
  A(\theta) \;=\; -\frac{2\nu}{\pi}\,\sec^{3}\theta
    \sum_{j} \overline{F_j(\sec\theta)}\;
      \exp\!\bigl[-\ii\,\nu\sec\theta\,(x_j + y_j\tan\theta)\bigr],
  \;}
  \label{eq:Adef}
\end{equation}
with $\overline{(\cdot)}$ the complex conjugate and $(x_j, y_j)$ hull $j$'s
midpoint in fleet coordinates (\texttt{amp\_pair}). The
conjugate and the $-\tfrac{2\nu}{\pi}\sec^{3}\theta$ prefactor together fix both
magnitude and phase: they follow Tuck, Scullen \& Lazauskas (2001, 2002) mapped to this crate's
conventions, and reproduces the classical Havelock point-source phase (a bow
source's transverse system trails a crest first). Its magnitude is pinned by the
free-wave resistance identity, which reproduces \eqref{eq:michell-theta} term for
term:
\begin{equation}
  \Rw \;=\; \tfrac12\,\pi\rho U^{2}
    \int_{-\pi/2}^{\pi/2} |A(\theta)|^{2}\cos^{3}\theta \dd\theta .
  \label{eq:Rspec}
\end{equation}
The integrand of \eqref{eq:Rspec}, namely
$\tfrac12\pi\rho U^{2}|A(\theta)|^{2}\cos^{3}\theta$, is the angular
resistance density $\dd\Rw/\dd\theta$ (\texttt{resistance\_density}).

\subsection{Elevation grid}

The surface elevation is the real part of the inverse angular transform
\begin{equation}
  \zeta(x, y) \;=\; \operatorname{Re}
    \int_{-\pi/2}^{\pi/2} A(\theta)\,
      \exp\!\bigl[\ii\,\nu\sec\theta\,(x + y\tan\theta)\bigr]\dd\theta .
  \label{eq:zeta}
\end{equation}
This is sampled on an inclusive rectangular grid of $n_x \times n_y$ points
spanning $[x_0, x_1] \times [y_0, y_1]$, stored row-major over $y$
(\texttt{WaveGrid.zeta[iy*nx + ix]}); a longitudinal or transverse cut is the
$n_y = 1$ or $n_x = 1$ special case. The evaluation exploits three structures:

\begin{enumerate}
  \item \textbf{$\pm\theta$ pairing.} $F_j$ depends only on $\sec\theta$, so the
        two signed systems $A(+\theta)$ and $A(-\theta)$ share their inner-integral
        evaluations; only the transverse placement phase $e^{\pm\ii k_y y_j}$
        flips. For each $|\theta|$ the row contribution is
        $A(+\theta)e^{\ii k_y y} + A(-\theta)e^{-\ii k_y y}$.
  \item \textbf{Phasor recurrence in $x$.} Along a grid row the longitudinal
        phase $e^{\ii k_x x}$ (with $k_x = \nu\sec\theta$) advances by the constant
        factor $e^{\ii k_x h_x}$ per column, so each column phasor is obtained
        from the previous by one complex multiply.
  \item \textbf{Same panel marcher as the resistance.} The $\theta$-integral
        \eqref{eq:zeta} is integrated with oscillation-rate-sized Gauss--Legendre
        panels and truncated once the amplitude has decayed
        (Section~\ref{sec:outer}).
\end{enumerate}

\paragraph{Resolution and spectral limits.} Two cutoffs keep the grid physical
and alias-free:
\begin{itemize}
  \item \emph{Resolution cutoff.} A grid of spacing $(h_x, h_y)$ can only render
        components with $\nu\sec\theta\,h_x < \pi$ and
        $\nu\sec\theta\tan\theta\,h_y < \pi$ (at least $\sim$2 points per
        wavelength). The largest resolvable $\theta$ is found by bisection, and
        the last $30\%$ (in $\sec^2\theta$) before it is cosine-tapered to zero to
        suppress aliasing. If even the transverse wavelength $2\pi/\nu$ is
        unresolved, the grid is rejected. Grids limited by this cutoff report
        \texttt{resolution\_limited = true}.
  \item \emph{Spectral cap.} Components beyond $\sec\theta = 15$ (wavelength below
        $\lambda_{\mathrm{tr}}/225$) are always dropped: for wall-sided hulls
        $|A|$ decays only algebraically toward the track, so the diverging system
        carries ever-shorter ripple that viscosity would destroy in reality (cf.\
        the eddy-viscosity factor of Tuck, Scullen \& Lazauskas 2002).
\end{itemize}

\paragraph{Validity.} Equation~\eqref{eq:zeta} is the far-field \emph{free-wave}
part of the linear solution. It omits the local (non-radiating) near-field
disturbance, so it is the complete linear answer only aft of the stern; on,
abreast of, or ahead of the hull it is not physical. Being linear, it also
cannot break: steep bow systems are indicative only.

\section{Data flow summary}
\label{sec:dataflow}

\begin{center}
\begin{tabular}{r c l}
  \toprule
  \multicolumn{3}{c}{\textbf{Resistance}} \\
  \midrule
  control net $P_{ij}$, knots & $\longrightarrow$ & B-spline surface $f$ (Eq.~\ref{eq:surface}) \\
  per span: corner partials $D_{ab}$ & $\longrightarrow$ & slope coeffs $c_{ab}$ (Eq.~\ref{eq:fxcoeff}, \texttt{fx\_coeff}) \\
  $c_{ab}$, moments $M_a, N_b$ & $\longrightarrow$ & amplitude $F(\lambda) = I + \ii J$ (Eq.~\ref{eq:Fexact}, exact) \\
  $|F|^2$ (or $|A|^2$ for fleets) & $\longrightarrow$ & outer $\theta$-quadrature (Eq.~\ref{eq:michell-theta}) \\
  & $\longrightarrow$ & wave resistance $\Rw$ \\
  \midrule
  \multicolumn{3}{c}{\textbf{Wave pattern}} \\
  \midrule
  $F_j(\sec\theta)$, placements & $\longrightarrow$ & amplitude density $A(\theta)$ (Eq.~\ref{eq:Adef}) \\
  $A(\theta)$, grid $(x,y)$ & $\longrightarrow$ & inverse transform $\zeta(x,y)$ (Eq.~\ref{eq:zeta}) \\
  & $\longrightarrow$ & elevation grid \texttt{WaveGrid} \\
  \bottomrule
\end{tabular}
\end{center}

\noindent The two branches share the amplitude functions $F_j$: the resistance
integrates $|F|^2\sec^3\theta$ over angle, while the wave pattern inverse-transforms
$A \propto \overline{F}\sec^3\theta$ back to physical space. Both are exact in the
geometry; the only approximation is the one-dimensional angular quadrature.

\end{document}
